\documentclass[12pt,a4paper]{article}

% ---------------------------
% Packages
% ---------------------------
\usepackage[margin=1in]{geometry}
\usepackage{amsmath,amssymb}
\usepackage{siunitx}
\usepackage{graphicx}
\usepackage{booktabs}
\usepackage{authblk}
\usepackage{hyperref}
\usepackage{url}
\usepackage{enumitem}
\usepackage{caption}
\usepackage{subcaption}
\usepackage{float}
\usepackage{listings}
\usepackage{xcolor}
\usepackage{microtype}
\usepackage{newtxtext,newtxmath}  % high-quality Times-like fonts

\usepackage{draftwatermark}
\SetWatermarkText{UNPUBLISHED CONFIDENTIAL}
\SetWatermarkScale{0.4}      % size of the text
\SetWatermarkColor[gray]{0.85} % light gray
\SetWatermarkAngle{55}       % rotation angle

\graphicspath{{figures/}}

% ---------------------------
% Listings setup (for code)
% ---------------------------
\lstdefinestyle{code}{
  basicstyle=\ttfamily\small,
  keywordstyle=\color{blue},
  commentstyle=\color{gray},
  stringstyle=\color{orange},
  showstringspaces=false,
  breaklines=true,
  frame=single
}

\lstset{style=code}

% ---------------------------
% Title / Authors
% ---------------------------
\title{LunarAtlas: Cleaning and Structuring Chandrayaan-3 LIBS Level-1 Data for Reproducible Lunar Spectroscopy}

\author[1]{Lovekesh Anand}
\author[1]{Dua Saeed}

\affil[1]{Independent Researchers, India}

\date{April 2026}

% ---------------------------
% Document
% ---------------------------
\begin{document}
\maketitle

\begin{abstract}
We present LunarAtlas, a backend-centric scientific data infrastructure that converts publicly available Chandrayaan-3 Laser-Induced Breakdown Spectroscopy (LIBS) Level-1 products into cleaned, long-form, and machine-accessible spectra suitable for quantitative lunar science. Starting from PDS4-compliant calibrated Level-1 tables, LunarAtlas implements a reproducible Python pipeline that reshapes wide Chandrayaan-3 LIBS files into per-wavelength records, pairs plasma and background measurements using mission metadata, and performs physically motivated background subtraction to produce cleaned spectra with preserved emission-line features.

On top of this cleaning pipeline, LunarAtlas provides a PDS4-aware PostgreSQL schema that stores spectral samples together with complete observation and instrument context, MD5 digests for all ingested files, and algorithm-version metadata for every processing step. We further introduce an adaptive min--max downsampling strategy for interactive visualization of high-resolution spectra, together with an analytic zoom-saturation analysis that defines where downsampling is meaningful and where the system must fall back to raw data.

Using processed Chandrayaan-3 LIBS spectra, we demonstrate the effect of cleaning on real Level-1 data (baseline suppression factor \(\sim 7\) in representative windows), quantify noise behaviour, and show how LunarAtlas exposes cleaned, NIST-validated spectra and provenance through a versioned API. This work establishes a practical blueprint for turning PDS4 Level-1 lunar LIBS products into a reusable, scientifically rigorous data resource.
\end{abstract}

\noindent\textbf{Keywords:} Laser-Induced Breakdown Spectroscopy; Chandrayaan-3; Lunar Regolith; Planetary Data Systems; PDS4; Data Cleaning; Spectral Visualization; Peak Preservation; NIST Validation.

% =================================================
\section{Introduction}
\label{sec:intro}
% =================================================

\subsection{Background and Motivation}

Laser-Induced Breakdown Spectroscopy (LIBS) is a versatile technique for \emph{in situ} elemental analysis in planetary environments.\cite{Cremers2006} By focusing a pulsed laser on a target, a transient plasma plume is created whose characteristic emission lines encode the elemental composition of the ablated material. Since its first planetary deployment on NASA's Mars Science Laboratory (MSL) rover \emph{Curiosity} in 2012, LIBS has become central to Martian geochemistry via the ChemCam instrument.\cite{Wiens2012,Maurice2012,Mangold2016}

The ChemCam mast unit combines a laser and telescope with body-mounted spectrometers spanning roughly 240--905~nm, and has acquired more than 188\,000 spectra from thousands of targets in Gale Crater.\cite{Mangold2016} SuperCam on the Mars 2020 \emph{Perseverance} rover extends this capability with co-aligned LIBS, Raman, and VISIR spectroscopy.\cite{Wiens2020,Bernardi2021} China's Tianwen-1 mission carries the Mars Surface Composition Detector (MarSCoDe), which combines LIBS and short-wave infrared reflectance spectroscopy on the Zhurong rover.\cite{Xu2021,MarSCoDe2024}

India's Chandrayaan-3 mission, which successfully landed near the lunar south pole in August 2023, marks the first deployment of LIBS on the Moon.\cite{ISROCh3web} The Pragyan rover carries a LIBS instrument for elemental composition analysis and an Alpha Particle X-ray Spectrometer (APXS) for complementary geochemistry.\cite{ISROCh3web} Operating in the near-vacuum lunar environment introduces different plasma dynamics and calibration constraints compared to Mars, making careful handling of Level-1 products essential.

\subsection{The Challenge of Using Chandrayaan-3 LIBS Level-1 Data}

Chandrayaan-3 LIBS data are distributed as calibrated Level-1 (L1) products conforming to Planetary Data System version~4 (PDS4) standards.\cite{Hughes2018} In the dataset used here, representative L1 files contain:

\begin{itemize}[nosep]
  \item an XML label describing mission/instrument context, observation times, and table structures;
  \item a CSV or fixed-width table where the first columns store metadata (e.g., \texttt{Time}, \texttt{Measurement Count}, \texttt{Operation Mode}, \texttt{Measurement Type}, \texttt{Force Reset Status}, \texttt{Laser Fire Status}), followed by hundreds to thousands of columns whose \emph{names} are wavelength values in nanometers (e.g., \texttt{164.35}, \texttt{164.74}, \ldots, \texttt{878.26}).
\end{itemize}

For one of the L1 files used in this work (\texttt{ch3\_lib\_002\_20230825T104221\_00\_l1.csv}), there are 2049 wavelength channels spanning 164.35--878.26~nm. Across mission products, the total number of wavelength samples per spectrum typically lies in the \(\sim 2\times 10^3\) to \(\sim 3\times 10^4\) range, depending on instrument configuration and processing level.

Although these products are already \emph{calibrated} at the Level-1 stage, they are not immediately ready for analysis:

\begin{enumerate}[label=(\roman*),nosep]
  \item They are in a wide format that is inconvenient for most scientific and database workflows.
  \item Plasma shots and background shots are mixed in the same tables and must be correctly identified and paired.
  \item Background subtraction and quality filtering must be performed consistently to obtain cleaned spectra.
  \item There is no out-of-the-box infrastructure for storing cleaned spectra with full provenance or serving them to clients.
\end{enumerate}

For Chandrayaan-3, the practical bottleneck is therefore not “calibrating” the instrument, but \emph{cleaning and structuring already calibrated Level-1 data} so that they can be used reproducibly.

\subsection{Research Objectives and Contributions}

The central question addressed by LunarAtlas is:

\medskip
\noindent\emph{How can Chandrayaan-3 LIBS Level-1 data be cleaned and structured into a reproducible, queryable, and machine-accessible resource without compromising scientific validity or peak fidelity?}
\medskip

This paper makes the following contributions, ordered by importance:

\begin{enumerate}[label=(C\arabic*),leftmargin=2em]
  \item \textbf{LIBS Level-1 cleaning pipeline:} A reproducible Python pipeline that ingests Chandrayaan-3 LIBS Level-1 tables, reshapes wide PDS4 products into long-form spectra, identifies and pairs plasma/background shots using mission metadata, and performs physically motivated background subtraction to produce cleaned spectra while preserving emission-line features.
  \item \textbf{PDS4-aware data model:} A normalized PostgreSQL schema that stores cleaned spectral samples together with mission, instrument, observation, and measurement context, enabling SQL-level queries over wavelengths, observation geometry, and processing versions while preserving PDS4 provenance and MD5 checksums.
  \item \textbf{Integrity and version tracking:} For every source and derived file, LunarAtlas records MD5 digests, file sizes, ingestion timestamps, and processing algorithm versions, enabling bit-level integrity checks and traceability from any visualization back to exact Level-1 products and pipeline configurations.
  \item \textbf{Adaptive min--max downsampling and zoom analysis:} A mathematically defined min--max downsampling method for spectra, together with an analytic zoom-saturation criterion that identifies the zoom range where downsampling is informative and the range where the system must fall back to raw data.
  \item \textbf{API and visualization layer:} A backend-centric API and lightweight frontend that expose cleaned spectra, NIST-based peak matches, and provenance in a way that can be reused for future lunar and planetary LIBS missions.
\end{enumerate}

The remainder of the paper is organized as follows. Section~\ref{sec:related} reviews related work. Section~\ref{sec:data} describes the Chandrayaan-3 LIBS dataset and PDS4 context. Section~\ref{sec:architecture} presents the LunarAtlas architecture and data model. Section~\ref{sec:processing} details the Level-1 cleaning pipeline. Section~\ref{sec:math} covers the adaptive downsampling formulation and zoom analysis. Section~\ref{sec:api} describes the API and visualization workflow. Section~\ref{sec:results} presents cleaning and zoom-range results. Section~\ref{sec:discussion} discusses implications and limitations. Section~\ref{sec:conclusion} concludes.

% =================================================
\section{Related Work}
\label{sec:related}
% =================================================

\subsection{Planetary LIBS Instruments and Mission Heritage}

ChemCam on MSL provides remote compositional information using the first planetary LIBS instrument, accompanied by a remote micro-imager.\cite{Wiens2012,Maurice2012,Mangold2016} SuperCam on Mars 2020 extends this approach with co-aligned LIBS, Raman, VISIR, and acoustic measurements.\cite{Wiens2020,Bernardi2021} MarSCoDe on Tianwen-1 combines LIBS and SWIR spectroscopy to cover 240--2400~nm for Martian surface studies.\cite{Xu2021,MarSCoDe2024} These missions demonstrate LIBS maturity and motivate robust data infrastructures that can handle diverse instruments and calibration schemes.

\subsection{Planetary Data Systems and Archival Standards}

PDS4 organizes data into bundles, collections, and products, using XML labels to describe structure, units, and provenance.\cite{Hughes2018} For LIBS and other spectrometers, typical products pair XML labels with CSV or fixed-width tables containing wavelength, intensity, and flags (laser status, plasma/background, timing). While self-describing, these archives lack ready-to-use APIs or visualization tools, and users must write bespoke parsers or rely on libraries that expose PDS4 tables as generic raster or vector layers.\cite{GDALPDS4web}

\subsection{Downsampling and Feature-Preserving Visualization}

General-purpose time-series downsampling methods such as mean/median aggregation, Douglas--Peucker simplification, and Largest Triangle Three Buckets (LTTB) are widely used in web visualization libraries, including JavaScript charting frameworks.\cite{amChartsWeb} They prioritize perceived shape over strict retention of narrow peaks. For LIBS and other spectroscopy, however, narrow emission lines are the primary information carriers.

Spectral analysis literature focuses on baseline correction, peak detection, line fitting, and chemometric modelling, typically assuming access to full-resolution data.\cite{Cremers2006} The specific problem of interactive, peak-preserving visualization of high-resolution planetary LIBS spectra, starting from PDS4 Level-1 products, has not been comprehensively addressed.

% =================================================
\section{Chandrayaan-3 LIBS Dataset and PDS4 Context}
\label{sec:data}
% =================================================

\subsection{Mission Overview and Landing Site}

Chandrayaan-3 launched on July 14, 2023 and achieved a soft landing near the lunar south pole on August 23, 2023.\cite{ISROCh3web} The Pragyan rover carried a LIBS instrument for elemental composition analysis and an APXS for complementary X-ray spectroscopy.\cite{ISROCh3web} The landing site is of high interest for access to ancient highland material and potential volatile-rich regions.

\subsection{LIBS Level-1 Data Products}

Chandrayaan-3 LIBS data are released as calibrated Level-1 (L1) products conforming to PDS4, consisting of:

\begin{itemize}[nosep]
  \item XML labels describing observation time, mission/instrument context, and table column definitions;
  \item tabular files storing wavelength channels (column names), intensities (cell values), measurement identifiers, timestamps, and flags (laser status, plasma/background).
\end{itemize}

For the L1 file \texttt{ch3\_lib\_002\_20230825T104221\_00\_l1.csv}, there are 2049 wavelength columns spanning 164.35--878.26~nm. Multiple measurements per observation (background + plasma shots) are represented as rows sharing common wavelength columns but different metadata values.

Across the broader Chandrayaan-3 LIBS archive, Level-1 products typically contain on the order of \(2\times 10^3\) to \(2.7\times 10^4\) wavelength–intensity pairs per measurement, depending on acquisition mode and processing level.

\subsection{PDS4 Semantics and Database Mapping}

LunarAtlas maps Level-1 semantics into a relational schema with tables for mission, instrument, observation, file version, measurement, and spectral samples. PDS4 logical identifiers (LIDs), timestamps, and MD5 digests are preserved in a \texttt{file\_version} table, while long-form spectral samples are stored in a \texttt{spectral\_data} table keyed by measurement and wavelength. This mapping enables SQL-level queries on wavelength ranges, measurement types, and observation context while maintaining traceability back to original PDS4 products.

% =================================================
\section{LunarAtlas System Architecture}
\label{sec:architecture}
% =================================================

\subsection{Design Principles}

LunarAtlas is guided by three principles:

\begin{enumerate}[nosep]
  \item \textbf{Single source of scientific truth:} All scientific logic (cleaning, bucketing, peak detection, NIST matching, integrity checks) resides in the backend; the frontend is a rendering layer only.
  \item \textbf{Reproducibility by design:} Every ingestion and processing step is versioned; API responses include algorithm versions and, where relevant, commit hashes; each ingested file has an MD5 checksum.
  \item \textbf{Separation of concerns:} Data persistence, scientific logic, and presentation are implemented in distinct tiers.
\end{enumerate}

\subsection{Three-Tier Architecture and Integrity Tracking}

The architecture comprises:

\begin{itemize}[nosep]
  \item \textbf{Data layer:} PostgreSQL stores mission metadata, measurement tables, spectral samples, NIST reference lines, and MD5 digests for every source and derived file.
  \item \textbf{Logic layer:} A Python backend (FastAPI/Flask) implements PDS4 parsing, ingestion, Level-1 cleaning, bucket computation, min--max aggregation, peak detection, NIST matching, MD5 computation, and JSON serialization.\cite{FastAPIweb}
  \item \textbf{Presentation layer:} A React + amCharts frontend manages viewports and chart rendering, never performing scientific calculations or client-side downsampling.\cite{amChartsWeb}
\end{itemize}

For each ingested file, the \texttt{file\_version} table stores original file name and size, MD5 checksum, ingestion timestamp, and processing algorithm version. This allows any visualization or analysis to be traced back to an exact Level-1 product and pipeline configuration.

% =================================================
\section{Data Processing and Cleaning Pipeline}
\label{sec:processing}
% =================================================

\subsection{Overview of the Python Pipeline}

The cleaning pipeline is implemented in Python and consists of two main scripts:

\begin{itemize}[nosep]
  \item \texttt{batch\_process\_libs.py}: reads Chandrayaan-3 LIBS Level-1 CSV files, reshapes wide tables into long-form per-wavelength rows, associates each spectral sample with measurement-level metadata, performs background subtraction and quality filtering, computes MD5 checksums, and writes cleaned spectra to disk and/or the database.
  \item \texttt{plot\_libs\_spectra.py}: loads specific raw and cleaned spectra, generates diagnostic plots (full-range and zoomed) that compare Level-1 and cleaned spectra, and exports figures used in this paper.
\end{itemize}

This pipeline is the core contribution of LunarAtlas: it systematically transforms calibrated Level-1 data into cleaned, structured spectra that can be queried and visualized reproducibly.

\subsection{From Wide Level-1 Tables to Long-Form Spectra}

Chandrayaan-3 Level-1 LIBS products store intensities in a wide format:

\begin{itemize}[nosep]
  \item initial metadata columns: \texttt{Time}, \texttt{Measurement Count}, \texttt{Operation Mode}, \texttt{Measurement Type}, \texttt{Force Reset Status}, \texttt{Laser Fire Status}, etc.;
  \item wavelength-value columns: numeric column names interpreted as wavelength channels in nanometers (e.g., 164.35, 164.74, \ldots, 878.26).
\end{itemize}

The cleaning pipeline identifies all columns whose names can be parsed as floating-point wavelength values and reshapes each Level-1 row into long-form:

\[
\texttt{Measurement\_ID},\ \texttt{Time\_UTC},\ \texttt{Measurement\_Count},\ \texttt{Wavelength\_nm},\ \texttt{Raw\_Intensity},\ \ldots \tag{1}
\]

This representation is natural for both database storage and subsequent spectral operations (background pairing, peak detection, NIST matching).

\subsection{Background Subtraction and Plasma Selection}

For each observation sequence containing plasma and background shots, the pipeline:

\begin{enumerate}[nosep]
  \item identifies valid plasma measurements using metadata flags (e.g., \texttt{Laser Fire Status} and a derived boolean \texttt{Is\_Valid\_Plasma = True}, \texttt{Is\_Background = False});
  \item pairs each plasma spectrum with a corresponding background spectrum (matched in time and configuration);
  \item computes cleaned intensities via
  \[
  I_{\text{cleaned}}(\lambda) = \max\bigl(0,\, I_{\text{plasma}}(\lambda) - I_{\text{background}}(\lambda)\bigr); \tag{2}
  \]
  \item marks each cleaned sample with validity flags while retaining instrument metadata such as delay time, integration time, number of pulses, and laser energy.
\end{enumerate}

The resulting cleaned table has columns such as:

\[
\texttt{Measurement\_ID},\ \texttt{Time\_UTC},\ \texttt{Measurement\_Count},\ \texttt{Wavelength\_nm},\ \texttt{Cleaned\_Intensity},\ \texttt{Is\_Valid\_Plasma},\ \texttt{Is\_Background},\ \ldots \tag{3}
\]

and is written to disk and/or ingested into the \texttt{spectral\_data} table.

\subsection{Effect of Cleaning on a Chandrayaan-3 Spectrum}

We demonstrate the pipeline using one plasma spectrum from \texttt{ch3\_lib\_002\_20230825T104221\_00\_l1.csv} and its corresponding cleaned spectrum from \texttt{ch3\_lib\_002\_20230825T104221\_00\_l1\_cleaned.csv}. Wavelength channels span 164.35--878.26~nm.

\begin{figure}[H]
  \centering
  \includegraphics[width=\textwidth]{libs_raw_vs_cleaned_full.png}
  \caption{Full-range Chandrayaan-3 LIBS spectrum. Raw Level-1 counts (blue) and cleaned intensity (orange) for a single plasma measurement spanning 164.35--878.26~nm. Cleaning removes continuum background while preserving emission features.}
  \label{fig:raw_clean_full}
\end{figure}

Figure~\ref{fig:raw_clean_zoom} shows a zoomed region from 380 to 420~nm. Background subtraction flattens the baseline and enhances the contrast of emission lines relative to the local continuum.

\begin{figure}[H]
  \centering
  \includegraphics[width=\textwidth]{libs_raw_vs_cleaned_zoom.png}
  \caption{Zoomed spectrum (380--420~nm) for the same measurement as Fig.~\ref{fig:raw_clean_full}. Raw L1 counts (blue) and cleaned intensity (orange) are shown. Cleaning removes background structure while preserving line positions.}
  \label{fig:raw_clean_zoom}
\end{figure}

\subsection{File Integrity and MD5 Checksums}

At ingestion, the pipeline computes MD5 checksums for every source Level-1 file and every exported cleaned spectrum file. Recorded metadata include original file name and size, MD5 checksum (128-bit hex), ingestion timestamp, and processing algorithm version. These fields enable:

\begin{itemize}[nosep]
  \item verification that archived files have not been corrupted or modified,
  \item traceability from any API response or visualization back to a specific Level-1 product and cleaning version,
  \item future re-processing studies that compare outputs across algorithm versions.
\end{itemize}

% =================================================
\section{Mathematical Formulation of Adaptive Min--Max Downsampling}
\label{sec:math}
% =================================================

\subsection{Bucket Size as a Function of Zoom}

Given a viewing window
\[
S = \{(\lambda_i, I_i): i=1,\ldots,N_{\text{raw}}\},\qquad
\lambda_{\min} \le \lambda_i \le \lambda_{\max},
\]
define \(\Delta\lambda = \lambda_{\max} - \lambda_{\min}\). For zoom level \(z\in\mathbb{Z}_{\ge 0}\), the unconstrained bucket size is
\[
b_{\text{size}}(z) = \frac{\Delta\lambda}{\text{BASE\_BUCKETS}\times 2^z}. \tag{4}
\]
With \(\text{BASE\_BUCKETS}\approx 1000\), this gives roughly one bucket per pixel at \(z=0\).

To avoid buckets smaller than the instrument resolution, a minimum bucket size \(\text{MIN\_BUCKET\_SIZE}\) (e.g., \SI{0.01}{nm}) is enforced:
\[
b_{\text{size,final}} = \max\bigl(b_{\text{size}}(z), \text{MIN\_BUCKET\_SIZE}\bigr). \tag{5}
\]

\subsection{Zoom Saturation Threshold and Valid Zoom Range}

The unsaturated downsampling regime is defined by
\[
b_{\text{size}}(z) \ge \text{MIN\_BUCKET\_SIZE}
\quad\Longleftrightarrow\quad
\frac{\Delta\lambda}{\text{BASE\_BUCKETS}\times 2^z} \ge \text{MIN\_BUCKET\_SIZE}. \tag{6}
\]

The maximum unsaturated zoom level is
\[
z_{\max}(\Delta\lambda) = \left\lfloor
\log_2\left(\frac{\Delta\lambda}{\text{BASE\_BUCKETS}\times \text{MIN\_BUCKET\_SIZE}}\right)
\right\rfloor. \tag{7}
\]

For \(z > z_{\max}\), bucket size clamps to \(\text{MIN\_BUCKET\_SIZE}\) and the implementation switches to raw-data mode, returning all spectral samples in the requested window.

\subsection{Worked Example: Saturation in a 300--400~nm Window}

For window [300, 400]~nm (\(\Delta\lambda = 100\)~nm) with \(\text{BASE\_BUCKETS} = 1000\), \(\text{MIN\_BUCKET\_SIZE} = 0.01\)~nm:

\[
b_{\text{size}}(6) = \frac{100}{1000\times 2^6} = \frac{100}{64{,}000} = \frac{1}{640} \approx 1.56\times 10^{-3}\ \text{nm}, \tag{8}
\]
so applying the floor constraint yields
\[
b_{\text{size,final}} = \max\bigl(1.56\times10^{-3}, 0.01\bigr) = 0.01\ \text{nm},\qquad
N_{\text{buckets}} = \left\lceil \frac{100}{0.01} \right\rceil = 10{,}000. \tag{9}
\]

At \(z=6\) the algorithm has saturated: each bucket is at the minimum allowed width and further zoom increases cannot reduce bucket size. In this regime LunarAtlas uses raw-mode inspection.

\subsection{Representative Bucket Sizes and $z_{\max}$}

For the Chandrayaan-3 spectrum analysed in Section~\ref{sec:processing} (164.35--878.26~nm, \(\Delta\lambda = 713.91\)~nm) and \(\text{BASE\_BUCKETS} = 1000\), \(\text{MIN\_BUCKET\_SIZE} = 0.01\)~nm, we obtain:

\begin{table}[H]
  \centering
  \caption{Zoom-saturation analysis. Bucket sizes and maximum unsaturated zoom levels for representative wavelength windows. Beyond \(z_{\max}\), the system switches to raw-mode.}
  \label{tab:zmax}
  \begin{tabular}{@{}lccc@{}}
    \toprule
    Window & $\Delta\lambda$ (nm) & $z_{\max}$ & $b_{\text{size}}(z_{\max})$ (nm) \\
    \midrule
    Full Chandrayaan-3 band & 713.91 & 6 & 0.0112 \\
    100~nm window           & 100.00 & 3 & 0.0125 \\
    10~nm window            & 10.00  & 0 & 0.0100 \\
    \bottomrule
  \end{tabular}
\end{table}

These values determine how many logical zoom levels represent genuine downsampling versus raw-mode viewing.

% =================================================
\section{API Design and Visualization Workflow}
\label{sec:api}
% =================================================

\subsection{Spectral Endpoints}

The primary spectral endpoint is
\begin{verbatim}
GET /api/v1/spectral
\end{verbatim}
with query parameters such as \texttt{observation\_id}, \texttt{min\_wavelength}, \texttt{max\_wavelength}, \texttt{zoom\_level}, and \texttt{include\_nist}. Responses include metadata, algorithm version, raw point counts, bucket parameters, and arrays of bucket objects with optional NIST annotations. All responses are machine-readable JSON with explicit units and schema-consistent field names.

Additional endpoints expose mission/instrument metadata, measurement catalogs, MD5 digests, and NIST reference lines.

\subsection{Progressive Loading and Zoom Strategy}

Frontend interaction proceeds as follows:

\begin{enumerate}[nosep]
  \item request an overview spectrum (full range, low zoom) for a selected cleaned measurement;
  \item on zoom events, map viewport bounds to wavelength ranges and integer zoom levels;
  \item for each range, compute \(z_{\max}(\Delta\lambda)\) on the backend; if the requested zoom \(z\le z_{\max}\), return min--max buckets; otherwise, return raw spectral samples in the range;
  \item render buckets or raw data using amCharts, with optional NIST line overlays.
\end{enumerate}

This strategy ensures that high-zoom views are always scientifically faithful to the cleaned spectra and do not rely on synthetic sub-resolution buckets.

\subsection{Frontend Visualization}

The React + amCharts frontend renders spectral buckets as filled ranges between minimum and maximum intensity per bucket, plotted versus wavelength, with optional markers at NIST-matched peaks. Tooltips display wavelength/intensity ranges, point counts, and identified elements with confidence scores. All quantities are produced by the backend; the frontend performs no scientific computation.

% =================================================
\section{Results and Experimental Validation}
\label{sec:results}
% =================================================

\subsection{Baseline Suppression and Noise Behaviour}

To quantify the effect of cleaning on baseline level and noise in a relatively line-poor region, we analyse a 20~nm window from 560 to 580~nm for the same Chandrayaan-3 measurement used in Figs.~\ref{fig:raw_clean_full} and \ref{fig:raw_clean_zoom}. In this window, the raw Level-1 spectrum has a mean intensity of approximately 1005~counts with an RMS fluctuation about the mean of \(\sim 47\)~counts, while the cleaned spectrum has a mean intensity of \(\sim 141\)~counts and an RMS fluctuation of \(\sim 184\)~counts (computed directly from the Level-1 and cleaned CSV files).

These values imply:

\begin{itemize}[nosep]
  \item a baseline (mean) suppression factor of \(\sim 7.1\) between raw and cleaned spectra;
  \item an increase in high-frequency variance about the mean by a factor of \(\sim 3.9\).
\end{itemize}

The cleaning pipeline thus strongly reduces low-frequency background while enhancing local contrast. For line detection and NIST matching, the reduction in baseline level increases peak prominence even though local variance is larger.

\subsection{Zoom Range and Peak Preservation}

Using the zoom-saturation criterion in Section~\ref{sec:math}, LunarAtlas configures downsampling as follows for each wavelength window:

\begin{itemize}[nosep]
  \item for \(z \le z_{\max}(\Delta\lambda)\), min--max downsampling is applied with bucket sizes determined by Eq.~(4);
  \item for \(z > z_{\max}(\Delta\lambda)\), bucket size is clamped to \(\text{MIN\_BUCKET\_SIZE}\) and the API returns raw spectral samples.
\end{itemize}

Table~\ref{tab:zmax} shows that the full Chandrayaan-3 band allows unsaturated downsampling up to \(z_{\max}=6\), a 100~nm window up to \(z_{\max}=3\), and a 10~nm window is already saturated at \(z=0\). These ranges are used directly in the implementation to define the backend behaviour for each requested zoom.

\subsection{Performance Metrics}

On a modest server configuration, end-to-end response times (query + aggregation + serialization) remain below \(\sim \SI{500}{ms}\) for typical cleaned spectra across zoom levels, with payload sizes between \(\sim 10\) and \(\sim 400\)~kB under compression. Composite indices on \texttt{(observation\_id, wavelength)} provide substantial speedups for spectral-range queries, and min--max buckets reduce rendered points by one to two orders of magnitude compared to raw spectra while preserving peak amplitudes and positions.

% =================================================
\section{Discussion}
\label{sec:discussion}
% =================================================

LunarAtlas demonstrates that a backend-centric architecture combined with a dedicated Level-1 cleaning pipeline can bridge the gap between archival PDS4 products and interactive, scientifically meaningful lunar LIBS analysis. The primary contribution is the reproducible transformation of Chandrayaan-3 LIBS Level-1 tables into cleaned, long-form spectra with full provenance; the adaptive min--max downsampling and zoom analysis are designed around this cleaned data model.

By centralising cleaning, integrity checks, and peak-preserving downsampling in the backend, LunarAtlas ensures that visualizations are reproducible scientific views rather than client-side approximations. The same design pattern can be extended to other spectroscopic missions (Mars LIBS, Raman, XRF, imaging spectroscopy) provided that instrument-specific calibration and reference databases are integrated appropriately.

% =================================================
\section{Conclusion}
\label{sec:conclusion}
% =================================================

LunarAtlas shows how Chandrayaan-3 LIBS Level-1 data can be cleaned and structured into a reproducible, interactive, and scientifically rigorous spectral infrastructure. Key elements include a PDS4-aware ingestion pipeline, physically motivated plasma/background cleaning and quality checks, a mathematically defined adaptive min--max downsampling algorithm with explicit zoom-saturation analysis, and a backend-centric API that serves cleaned, NIST-validated spectra and provenance to thin visualization clients.

The system achieves strong baseline suppression, preserves emission-line features, provides well-understood zoom behaviour, and maintains sub-second response times for typical queries. As future lunar and planetary missions continue to produce large spectroscopic archives, approaches like LunarAtlas will be essential for turning calibrated Level-1 products into practically usable scientific resources.

% =================================================
\section*{Acknowledgments}
% =================================================

We acknowledge the Indian Space Research Organisation (ISRO) for the public release of Chandrayaan-3 LIBS data through PDS4-compliant archives, and the National Institute of Standards and Technology (NIST) for maintaining the Atomic Spectra Database that underpins spectral validation.\cite{ReaderNIST} We thank Aditya Bhardwaj (Scientist, Indian Space Research Organisation, National Remote Sensing Centre, and TEDxMAIT speaker) for public talks on India's space programme and real-time satellite data systems, which helped shape LunarAtlas's broader vision. We also acknowledge the open-source scientific Python ecosystem (NumPy, SciPy, Pandas, Astropy) and the planetary science community whose LIBS instrument and calibration work informs this study.

% =================================================
\section*{Author Contributions}
% =================================================

L.~Anand: Conceptualization, Level-1 cleaning pipeline design, system architecture, algorithm development, database design, API implementation, data analysis, writing---original draft and editing.

D.~Saeed: Data processing implementation, PDS4 parsing, NIST validation integration, visualization design, performance benchmarking, writing---review and editing.

% =================================================
\section*{Data Availability}
% =================================================

Processed Chandrayaan-3 LIBS data products and LunarAtlas source code will be made publicly available upon manuscript acceptance. Raw Chandrayaan-3 LIBS data are available from ISRO's PDS4 archives.

% =================================================
\section*{Competing Interests}
% =================================================

The authors declare no competing interests.

% =================================================
% References come from your existing .bib or manual list.
% =================================================

\end{document}